% ======================================================================
%  7. Conclusion
% ======================================================================
\section{Conclusion}
\label{sec:conclusion}

This internship was motivated by a methodological question raised by Marco
Bonizzato in the context of scientific applications of Bayesian optimization.
At the sciNeurotech Lab, Bayesian optimization is a natural tool for optimizing
expensive and unknown responses to neural stimulation. More broadly, the project
was driven by the need for a clear modelling principle that can prevent
task-specific choices from being incorporated into Bayesian-optimization models
without sufficient theoretical justification (A problem that can naturally arise in applied work, as encoding the task directly in the surrogate is often the most intuitive modelling choice.). The guiding idea is that, whenever
possible, the surrogate should model the underlying phenomenon, while
task-specific structure should be introduced at the decision or acquisition
level. This principle leads to a common structural question, investigated here
across target matching, multi-objective optimization and prior-guided
optimization: \emph{where should task-specific information enter the
Bayesian-optimization pipeline?}


The answer developed here is a separation-of-concerns principle:
\begin{center}
\emph{the phenomenon may remain the same while the scientific task changes; model
the phenomenon $f$ faithfully, and encode the task in the acquisition whenever
possible.}
\end{center}

For target matching, keeping the task out of the model preserves RKHS regularity
and, through a pushforward of the confidence interval, a GP-UCB regret guarantee
(Theorem~\ref{thm:tm-soc}); folding the target into the model forces the surrogate
to represent a cusp, at a quantifiable RKHS cost (Theorem~\ref{thm:tm-barrier}).
For multi-objective optimization, the same pushforward argument covers any
Lipschitz scalarization (Theorem~\ref{thm:mo-pushforward}), and separating the
task from the model is precisely what lets one vector observation inform every
future preference. For prior-guided optimization, a spatial prior can be placed in
the acquisition with a controlled, never-vanishing effect, whereas a functional
prior estimated from an adaptively chosen design can lock onto a wrong belief
through a self-confirming loop (Theorem~\ref{thm:prior-failure}). Three controlled
experiments support the same conclusion, most strongly where the transformed
geometry is deceptive.

What the internship establishes is therefore a coherent case for keeping two
levels separate - the statistical model of the latent phenomenon and the
task-specific decision transformation - together with several mechanisms
explaining why this separation can matter, including regularity mismatch,
confidence-set pushforward, and self-confirming feedback under adaptive sampling.
Several directions remain open. A first priority is to deepen the probabilistic
analysis of misspecification, beyond RKHS approximation, in order to understand
how directly modelling a transformed objective distorts posterior uncertainty and
calibration. A second is to broaden the empirical validation through controlled
ablations and diagnostic benchmarks designed to isolate these mechanisms. Finally,
extending the separation-of-concerns principle to constrained Bayesian
optimization would test its relevance in a setting where task information enters
through feasibility rather than through the objective itself. A genuine lower
regret bound for the task-in-model strategy would provide a stronger algorithmic
separation, but our current results should instead be viewed as approximation and
misspecification results rather than as evidence that such a failure is
unavoidable.
